%% kapa: exam=1 round=29 year=2018 period=1 form=A subject=civil_law subject_ko=민법 seq=30 number=30
%% source: Q-Net 감정평가사 기출문제 (cst00309) artlSeq=5202808 file=extracted/1008/5202808/2018년 제29회 감정평가사 1차시험 기출문제/2018년 제29회 감정평가사 1차시험 기출문제/1교시 A형.pdf
%% conversion: kordoc 4.24.51 -> tools/gen1.py (auto); review=proofread-sonnet; flags=none
\begin{question}{30}
甲 소유 게임기 X를 乙이 빌려서 사용하던 중, 乙은 이러한 사정을 과실 없이 알지 못하는 丙에게 X를 50만원에 평온ㆍ공연하게 매도하고 점유를 이전해 주었다. 이에 관한 설명으로 옳은 것은? (다툼이 있으면 판례에 따름)
\begin{choices}
  \item 점유에는 공신력이 없으므로 丙은 X의 소유권을 선의취득할 수 없다.
  \item 乙과 丙 간의 매매계약이 무효이더라도 丙은 X의 소유권을 선의취득할 수 있다.
  \item 丙이 점유개정으로 점유를 취득하였더라도 X의 소유권을 선의취득할 수 있다.
  \item 만약 乙의 점유보조자가 X를 절취하여 丙에게 매도하였더라도 丙은 X의 소유권을 선의취득할 수 있다.
  \item 만일 X가 게임기가 아니라 건물인 경우에도 丙은 소유권을 선의취득할 수 있다.
\end{choices}
\end{question}
\answer{4}
